\section{Separation guarantees and their interpretation}\label{sec:guarantees-overview}
The common cut family permits a direct comparison of sampling policies, but cut validity alone does not establish that a numerical solver will close an objective gap. \Cref{sec:theory} supplies complete assumptions and proofs; the consequences most relevant to the experiments are as follows.

On a common compact domain with finitely many continuously differentiable convex rows, bounded gradients give bounded cut normals. For global or selected-term rows at integer candidates with row violation $v>0$, ECP separates by $v$. With a fixed strict anchor and positive row-specific slack margin $\delta$, ESH's exact boundary tangent separates by at least $\delta v/(LD)$, where $L$ bounds the gradient and $D$ bounds the anchor--candidate distance. Thus, for either policy, a fixed positive residual threshold implies a positive separation distance. Compactness then allows only finitely many newly rejected candidates that satisfy all previously enforced cuts (\cref{lem:margin,thm:integral}). This argument neither ranks the policies' iteration counts nor guarantees finite closure of an arbitrarily small objective gap.

For a fractional hull block, the relevant residual is
\[
 r(\nu,\lambda)=\lambda[g(\nu/\lambda)]_+,
\]
with value zero at the inactive origin. If $U$ bounds $[g]_+$ on the box, omitting that block is safe at tolerance $\eps_p$ whenever $\lambda U\le\eps_p$. Under the same compactness and enforcement conditions this yields finite fixed-residual fractional separation. A fixed cutoff $\lambda\le\tau$ instead bounds omitted residuals by $\tau U$, which need not meet the intended tolerance (\cref{thm:fractional}). The experiments use fixed cutoffs. Their LP stalls and small numerical residuals therefore cannot be read as certification of the exact continuous hull.

A strict interior permits repair of an approximate point in each separate disjunction, with a distance controlled by its weighted residuals. Intersecting disjunctions with global constraints can destroy a linear objective-error bound: \cref{ex:intersection} gives an explicit counterexample. Vanishing residuals and vanishing valid outer-master suboptimality give convergence in value under a common compact-domain assumption, without such a rate (\cref{prop:repair,prop:value}). For a single tree, old lazy cuts must be re-enforced and auxiliary cut generation must also terminate; \cref{thm:single-tree} states the required oracle contract. The frozen callbacks do not certify every premise of that contract.

The diagnostic in \cref{sec:oracle} isolates another reason for comparing actual work. For a fixed anchor, a convex, smooth increasing re-expression of a convex row can preserve its exact ESH boundary halfspace while weakening ECP at an exterior point, under the qualifications in \cref{prop:invariance}. A scalar example reduces ECP to a Newton recurrence with an arbitrarily long initial transient, although it converges locally quadratically for each fixed representation. The executable diagnostic counts both values and gradients, including the bisection overhead. Centered ellipsoids favor ESH's geometric cut, but an off-center disk gives explicit opposite witnesses: neither policy generally dominates. These results explain possible differences in cut selection without predicting whole-solver runtime.
